\subsection{Proof of \cref{lem:all-solutions-of-C-are-in-the-core}}
\label{proof:all-solutions-of-C-are-in-the-core}

\repi*

    Let $\eps=2^{-n}, \eps_0=1.6\eps$.
    Let $\mathcal{D}=\settxt{(x,y,z)\in \Delta^2: y\in 0.5\pm 0.1, z\in 0.2\pm 0.1}$ denote the core region.
    We are going to show that $\vec{x}^{(1)}, \vec{x}^{(2)}, \vec{x}^{(3)}\in \mathcal{D}$.
    We prove this fact by contradiction. 
    W.l.o.g., we suppose that $\vec{x}^{(1)}$ violates this condition.
    Since $\normtxt{\vec{x}^{(i)}-\vec{x}^{(j)}}{\infty}\leq \eps$, all of them are inside the following subset:
    \begin{align*}
        \overline{\mathcal{D}} = \set{\vec{y}\in \Delta^2: y_2\notin [0.4+\eps,0.6-\eps] \text{ or } y_3\notin [0.1+\eps,0.3-\eps]} ~.
    \end{align*}
    According to \cref{alg:base-instance} and \cref{def:rect-sperner}, we have 
    \begin{align*}
        \forall \vec{y}\in \overline{\mathcal{D}}, \quad C(\vec{y}) = \begin{cases}
            \hfill 1 \hfill & \text{if $y_2\leq 0.5$ and $y_3\leq 0.1$,} 
            \\
            \hfill 2 \hfill & \text{if $y_2> 0.5$ and $y_3\leq 0.1$,}\\
            \hfill 1 \hfill & \text{if $y_2\in [0.4,0.4+\eps_0)$ and $y_3\in [0.1, 0.1+\eps_0)$, }\\
            \hfill 2 \hfill & \text{if $y_2\in [0.6-\eps_0,0.6)$ and $y_3\in [0.1, 0.1+\eps_0)$, } \\
            1\text{ or }2 & \text{if $y_2\in [0.4+\eps_0, 0.6-\eps_0)$ and $y_3\in [0.1, 0.1+\eps]$,} \\
            \hfill 3 \hfill & \text{otherwise.}
        \end{cases}
    \end{align*}
    Next, we discuss several cases on where $\vec{x}^{(1)}\notin D$ is to establish the fact.
    \begin{enumerate}[itemsep=0em,topsep=0.5em]
        \item If $x^{(1)}_2\leq 0.4$, we have $C(\vec{x}^{(1)})\in \{1,3\}$. However, as any $\vec{y}\in \overline{\mathcal{D}}$ such that $C(\vec{y})=2$ has $y_2\geq 0.4+\eps_0>0.4+\eps$, the triangle formed by $\vec{x}^{(1)},\vec{x}^{(2)},\vec{x}^{(3)}$ cannot be trichromatic. 
        \item If $x^{(1)}_2\geq 0.6$, we have $C(\vec{x}^{(1)})\in \{2,3\}$. However, as any $\vec{y}\in \overline{\mathcal{D}}$ such that $C(\vec{y})=1$ has $y_2<0.6-\eps_0<0.6-\eps$, the triangle formed by $\vec{x}^{(1)},\vec{x}^{(2)},\vec{x}^{(3)}$ cannot be trichromatic. 
        \item If $x^{(1)}_3 \leq 0.1$, we have $C(\vec{x}^{(1)})\in \{1,2\}$. Note that any $\vec{y}\in \overline{\mathcal{D}}$ such that $C(\vec{y})=3$ and $y_3\leq 0.1+\eps$ has $y_2\notin [0.4,0.6]$. If the triangle formed by $\vec{x}^{(1)}, \vec{x}^{(2)}, \vec{x}^{(3)}$ is trichromatic, we should further have $x_2^{(1)}, x_2^{(2)}, x_2^{(3)}\leq 0.4+\eps$ or have $x_2^{(1)}, x_2^{(2)}, x_2^{(3)}\geq 0.6-\eps$. However, as discussed above, we cannot find $C(\vec{x}^{(j)})=2$ with the first condition and we cannot find $C(\vec{x}^{(j)})=1$ with the second condition. The triangle formed by $\vec{x}^{(1)},\vec{x}^{(2)},\vec{x}^{(3)}$ cannot be trichromatic. 
        \item If $x_3^{(1)}\geq 0.3$, we have $C(\vec{x}^{(1)})=3$. Note that there is no $\vec{y}\in \overline{\mathcal{D}}$ such that $y_3>0.2$. The triangle formed by $\vec{x}^{(1)},\vec{x}^{(2)},\vec{x}^{(3)}$ cannot be trichromatic. 
    \end{enumerate}


\subsection{Proof of \cref{lem:lipschitz-of-the-converter-on-boundaries}}
\label{proof:lipschitz-of-the-converter-on-boundaries}

\repii*

According to the characterization \cref{lem:symmetry-on-converted-coordinates}, for any $z\in [0,1]$, we have
\begin{align*}
    \rel(1-z,0,z),~\rel(0,1-z,z) ~\in~
    \begin{cases}
        \hfill \{0,1\} \hfill & \text{if $z\leq 0.1-2\eps^2$~;} \\
        \hfill \settxt{0} \hfill & \text{if $z\in (0.1-2\eps^2, 0.1-\eps^2]$~;} \\
        \hfill \set{0.5\eps^{-2}\cdot(z-0.1+\eps^2)} \hfill & \text{if $z\in 0.1\pm \eps^2$~;} \\
        \hfill \settxt{1} \hfill & \text{if $z\in [0.1+\eps^2, 0.1+2\eps^2)$~;} \\
        \hfill \{0,1\} \hfill & \text{if $z\geq 0.1+2\eps^2$~.}
    \end{cases}
\end{align*}
That is, when $z\in 0.1\pm 2\eps^2$, we have an exact characterization for it.
Note that the RHS of the characterization is continuous and $0.5\eps^{-2}\cdot (z-0.1+\eps^2)$ is $0.5\eps^{-2}$-Lipschitz. 
We have obtained this lemma for any $z,z'\in 0.1\pm 2\eps^{2}$.

On the other hand, w.l.o.g., suppose that $z\notin 0.1 \pm 2\eps^{2}$.
If $|z-z'|\geq \eps^2$, this lemma is simply trivial because we have $\eps^{-2}\cdot |z-z'|\geq 1$. 
Otherwise, according to the triangle inequality, $z'\notin 0.1\pm \eps^2$. 
Therefore, we have $\rel(\vec{x}), \rel(\vec{x'}) \in \{0,1\}$, which matches the second bullet of this lemma.

\subsection{Proof of \cref{fact:new-coordinate-converter-at-the-bottom-part}}
\label{proof:new-coordinate-converter-at-the-bottom-part}

\repiii*

Note that in our base instance $C$, we have 
\begin{align*}
    \forall \vec{x'}\in \Delta^2, \quad
    \begin{cases}
    C(\vec{y'}) = 1 & \quad \text{if } y'_3\leq 0.1, y'_2 \leq 0.5~,\\
    C(\vec{y'}) = 2 & \quad \text{if } y'_3\leq 0.1, y'_2 > 0.5~.
    \end{cases}
\end{align*}

If $y_2\leq 0.5$, we have $C(\vec{y})=1$. 
$C(\vec{y'})$ is different with $C(\vec{y})$ only if $y'_2>0.5$ or $y'_3>0.1$. 
Since $0.1-y_3>0.04\gg 2\eps^2$, $\vec{y}$ is hot or warm if and only if $\alpha(y_3) \cdot (0.5-y_2)<2\eps^2$, i.e., $g(\vec{y})\in (-2\eps^2,0)$. 
In this case, we have $d^{\alpha}(\vec{y}, \nna(\vec{y})) = \alpha(y_3)\cdot |0.5-y_2| = |g(\vec{y})|$, $C^{\alpha}_{\sf nn}(\vec{y})=2$, and 
\begin{align*}
    \rela(\vec{y}) = \left(0.5-0.5\eps^{-2}\cdot |g(\vec{y})|\right)_+ = \left(0.5+0.5\eps^{-2}\cdot g(\vec{y})\right)_{[0,1]}~.
\end{align*}

On the other hand, if $y_2> 0.5$, we have $C(\vec{y})=2$. 
$C(\vec{y'})$ is different with $C(\vec{y})$ only if $y'_2\leq 0.5$ or $y'_3>0.1$. 
Since $0.1-y_3>0.04\gg 2\eps^2$, $\vec{y}$ is hot or warm if and only if $\alpha(y_3) \cdot (y_2-0.5)<2\eps^2$, i.e., $g(\vec{y})\in (0,2\eps^2)$. 
In this case, we have $d^{\alpha}(\vec{y}, \nna(\vec{y})) = \alpha(y_3)\cdot |0.5-y_2| = |g(\vec{y})|$, $C^{\alpha}_{\sf nn}(\vec{y})=1$, and 
\begin{align*}
    \rela(\vec{y}) = \left(0.5+0.5\eps^{-2}\cdot |g(\vec{y})|\right)_- = \left(0.5+0.5\eps^{-2}\cdot g(\vec{y})\right)_{[0,1]}~.
\end{align*}

\subsection{Proof of \cref{lem:poly-time-recovery-symmetry}}
\label{proof:poly-time-recovery-symmetry}

\repiv*

We assume that $\vec{x}^{(1)}, \vec{x}^{(2)}, \vec{x}^{(3)}$ give a trichromatic triangle in the \outofk Approximate Symmetric \Sperner instance $\Csym{k}$ such that $\normtxt{\vec{x}^{(i)}-\vec{x}^{(j)}}\infty\leq 2^{-4kn}$, where $k\geq 3$.
As in the earlier proof of \cref{lem:poly-time-recovery}, 
we use $\vec{y}^{(i)}(\vec{x}) = (y^{(i)}_1(\vec{x}), y^{(i)}_2(\vec{x}), y^{(i)}_3(\vec{x}))$
to denote the {\em intermediate projections} and use
$\vec{c}^{(i)}(\vec{x}) = (c^{(i)}_1(\vec{x}), c^{(i)}_2(\vec{x}), c^{(i)}_3(\vec{x}))$ 
to denote the {\em intermediate palettes}.
In addition, we introduce a new notation $\tilde{y}_0(\vec{x})$ for the converted coordinate we compute on \cref{line:first-converted-coordinate}.
In the rest of this section, because the subscripts we will use for $\vec{c}^{(i)}$ can be very complicated, we will use $c_j^{(i)}(\vec{x})$ and $c^{(i)}(\vec{x}, j)$ interchangeably for better presentation.
For any vector $\vec{y}\in \Delta^2$, we use $i^*(\vec{y})$ to denote the first non-zero index of $\vec{y}$, i.e., 
\begin{align}
    \label{eqn:first-nonzero-index-symmetry}
    i^*(\vec{y}) = \begin{cases}
        1 & \text{if $y_1>0$,}\\
        2 & \text{if $y_1=0$ and $y_2>0$,}\\
        3 & \text{otherwise.}
    \end{cases}
\end{align}

Our recovery algorithm (formalized in \cref{alg:recover-2d-sol-symmetry}) simulates \cref{alg:approx-symmetric-sperner-kgeq3} for each $\vec{x}^{(j)}$.
At time $i\in [2,k-1]$ when we have computed $\vec{y}^{(i)}(\vec{x}^{(j)})$ for each $j\in [3]$, we examine whether one of them is inside a trichromatic region, i.e., whether $C$ has three different colors in $\mathcal{N}(\vec{y}^{(i)}(\vec{x}^{(j)}))$ for some $j\in [3]$.
According to our earlier discussions, this examination can be done in polynomial time.
After we finish the simulation, we simply output $\vec{y}^{(k)}(\vec{x}^{(1)}), \vec{y}^{(k)}(\vec{x}^{(2)}), \vec{y}^{(k)}(\vec{x}^{(3)})$ as a solution for $C$. 


It is clear that any output during the simulation phase of this recovery algorithm gives a valid solution for the 2{\rm D}-\Sperner instance $C$. 
To prove \cref{lem:poly-time-recovery-symmetry}, we only need to show that the three intermediate projections after the simulation form a valid solution for $C$ if we output after the simulation phase. 
Or equivalently, if each $\vec{y}^{(i)}(\vec{x}^{(j)})$ does not lie in a trichromatic region, the final converted coordinates, $\vec{y}^{(k)}(\vec{x}^{(1)}), \vec{y}^{(k)}(\vec{x}^{(2)}), \vec{y}^{(k)}(\vec{x}^{(3)})$, give a solution for $C$.

As in \cref{lem:poly-time-recovery}, our proof will be considers two cases; the first is where
the given solution $\vec{x}^{(1)}, \vec{x}^{(2)}, \vec{x}^{(3)}$  for $\Csym{k}$ further enjoys the following property: 
\begin{align}
\label{eqn:assumption-for-simple-proof-of-recovery-symmetry}
\forall 2\leq i\leq k, \forall j\in [3], \quad P_{i+1}^{(k-i)}(\vec{x}^{(j)})\leq 0.9~.
\end{align}
One benefit of first considering this case is that we don't have too crazy projections in \cref{alg:recover-2d-sol-symmetry}, which can help us significantly simplify the proof.
The formal intermediate technical result is presented in \cref{cor:no-sols-in-sim-implies-final-ones-symmetry}.
Later, in the second step, we will explain how to prove \cref{lem:poly-time-recovery-symmetry} in the  case where this property is not satisfied.

\paragraph{Case 1: the output satisfies \cref{eqn:assumption-for-simple-proof-of-recovery-symmetry}.}
We can continue to use the Lipschitzness of the projection step (\cref{lem:lipschitz-of-projection}) to obtain Lipschitzness for the intermediate projections.

\begin{lemma}
\label{lem:intermediate-projections-are-close-to-each-other-symmetry}
    Consider any $k\geq 2$ and any $\vec{x}^{(1)}, \vec{x}^{(2)}\in \Delta^k$.
    Suppose that $P_{i+1}^{(k-i)}(\vec{x}^{(j)})\leq 0.9$ for any $2\leq i\leq k$ and any $j\in [2]$.
    Also, suppose that $\mathcal{N}(\vec{y}^{(i)}(\vec{x}^{(j)}))$ is at most bichromatic for any $2\leq i<k$ and any $j\in [2]$.
    Then, for any $2\leq i\leq k$, we have Lipschitzness for the third coordinate of the intermediate projections:
    \begin{align*}
        \abstxt{y^{(i)}_3(\vec{x}^{(1)})-y^{(i)}_3(\vec{x}^{(2)})}\leq 2^{O(k)}\cdot \normtxt{\vec{x}^{(1)}-\vec{x}^{(2)}}\infty~.
    \end{align*}
    Furthermore, at least one of the following is satisfied for the second coordinates of the intermediate projections:
    \begin{itemize}[itemsep=0.1em, topsep=0.4em]
        \item {\bf Lipschitz:} $\abstxt{y_2^{(i)}(\vec{x}^{(1)})-y_2^{(i)}(\vec{x}^{(2)})}\leq 2^{3in+O(k)}\cdot \normtxt{\vec{x}^{(1)}-\vec{x}^{(2)}}\infty$, or
        \item {\bf both are on the left/right boundaries:} for any $j\in [2]$, we have either $y_1^{(i)}(\vec{x}^{(j)})=0$ or $y_2^{(i)}(\vec{x}^{(j)})=0$.
    \end{itemize}
\end{lemma}
\begin{proof}
    We prove this lemma by induction on $i$. The base case is when $i=2$. 
    We have $y^{(2)}_3(\vec{x})=P_3^{(k-2)}(\vec{x})$ and thus we can obtain the Lipschitzness for $y^{(2)}_3(\cdot)$ by
    \begin{align*}
        \abs{y_3^{(i)}(\vec{x}^{(1)})-y_3^{(i)}(\vec{x}^{(2)})}
        &\leq 
        \normtxt{\vec{P}^{(k-i)}(\vec{x}^{(1)})-\vec{P}^{(k-i)}(\vec{x}^{(2)})}{\infty}
        \\
        &\leq 
        110\cdot \normtxt{\vec{P}^{(k-i-1)}(\vec{x}^{(1)})-\vec{P}^{(k-i-1)}(\vec{x}^{(2)})}{\infty}
        \\
        &\leq
        \cdots
        \\
        &\leq 
        110^{k-i}\cdot \normtxt{\vec{x}^{(1)}-\vec{x}^{(2)}}{\infty}
        \leq 2^{7k}\cdot \normtxt{\vec{x}^{(1)}-\vec{x}^{(2)}}{\infty}~,
    \end{align*}
    which works for any $i\geq 2$. 
    At the same time, we can similarly obtain $\abstxt{P_2^{(k-1)}(\vec{x}^{(1)})-P_2^{(k-1)}(\vec{x}^{(2)})}\leq 2^{7k}\cdot \normtxt{\vec{x}^{(1)}-\vec{x}^{(2)}}\infty$. 
    Because the function $g(y_0)=(0.5+0.5\eps^{-2}\cdot (y_0-0.1))_{[0,1]}$ is clearly $0.5\eps^{-2}$-Lipschitz, we have
    \begin{align*}
        \abstxt{\tilde{y}_0(\vec{x}^{(1)}) - \tilde{y}_0(\vec{x}^{(2)})} \leq 2^{2n} \cdot \abstxt{P_2^{(k-1)}(\vec{x}^{(1)})-P_2^{(k-1)}(\vec{x}^{(2)})} \leq 2^{2n+7k} \cdot \normtxt{\vec{x}^{(1)}-\vec{x}^{(2)}}\infty~.
    \end{align*}
    Hence, we can the first bullet of the second statement of this lemma. 
    \begin{align*}
        \abstxt{y_2^{(2)}(\vec{x}^{(1)}) - y_2^{(2)}(\vec{x}^{(2)})} \leq \abs{\tilde{y}_0(\vec{x}^{(1)})-\tilde{y}_0(\vec{x}^{(2)})} + \abs{y^{(2)}_3(\vec{x}^{(1)})-y^{(2)}_3(\vec{x}^{(2)})} \leq 2^{2n+7k+1}\cdot \normtxt{\vec{x}^{(1)}-\vec{x}^{(2)}}\infty~.
    \end{align*}

    Consider any $i_0\geq 3$. 
    Suppose that we have proved this lemma for $i<i_0$. 
    Next, we consider when $i=i_0$.
    Note that we have $y_3^{(i)}(\vec{x})=P_{i+1}^{(k-i)}(\vec{x})$ for any $\vec{x}$.
    By the premise that $P_{i+1}^{(k-i)}(\vec{x}^{(j)})\leq 0.9$, we can use \cref{lem:lipschitz-of-projection} to get
    the first statement of this lemma by the earlier inequalities. 

    Next, we establish the second statement of this lemma, in which we need to prove at least one of the bullets is satisfied.
    For the induction hypothesis, we will use the following more specific version of the first bullet 
    \[
        \abstxt{y_2^{(i)}(\vec{x}^{(1)})-y_2^{(i)}(\vec{x}^{(2)})}\leq 2^{3in+7k+\log  2i}\cdot \normtxt{\vec{x}^{(1)}-\vec{x}^{(2)}}\infty~.
    \]
    Note that $y_2^{(i)}(\vec{x})=(1-y_3^{(i)}(\vec{x}))\cdot \rela(\vec{y}^{(i-1)}(\vec{x}))$. 
    Then, it is easy to obtain that
    \begin{align*}
        \abs{y_2^{(i)}(\vec{x}^{(1)})-y_2^{(i)}(\vec{x}^{(2)})} 
        &= 
        \abs{(1-y_3^{(i)}(\vec{x}^{(1)}))\cdot \rela(\vec{y}^{(i-1)}(\vec{x}^{(1)})) - (1-y_3^{(i)}(\vec{x}^{(2)}))\cdot \rela(\vec{y}^{(i-1)}(\vec{x}^{(2)}))}
        \\
        &\leq
        \rela(\vec{y}^{(i-1)}(\vec{x}^{(1)}))\cdot \abs{y_3^{(i)}(\vec{x}^{(1)})-y_3^{(i)}(\vec{x}^{(2)})} 
        \\
        & \qquad \quad + (1-y_3^{(i)}(\vec{x}^{(2)})) \cdot \abs{\rela(\vec{y}^{(i-1)}(\vec{x}^{(1)})) - \rela(\vec{y}^{(i-1)}(\vec{x}^{(2)})) }
        \\
        &\leq
        \abs{y_3^{(i)}(\vec{x}^{(1)})-y_3^{(i)}(\vec{x}^{(2)})} + \abs{\rela(\vec{y}^{(i-1)}(\vec{x}^{(1)})) - \rela(\vec{y}^{(i-1)}(\vec{x}^{(2)})) }
        \\
        &\leq 
        2^{7k}\cdot \normtxt{\vec{x}^{(1)}-\vec{x}^{(2)}}\infty + \abs{\rela(\vec{y}^{(i-1)}(\vec{x}^{(1)})) - \rela(\vec{y}^{(i-1)}(\vec{x}^{(2)})) }~.
    \end{align*}

    Suppose that our induction hypothesis gives the first bullet for $i-1$.
    Combining the Lipschitzness on the third coordinate, we have 
    \begin{align*}
        \normtxt{\vec{y}^{(i-1)}(\vec{x}^{(1)})-\vec{y}^{(i-1)}(\vec{x}^{(2)})}\infty &\leq (2^{3(i-1)n+7 k + \log 2(i-1)}+2^{7k}) \cdot \normtxt{\vec{x}^{(1)}-\vec{x}^{(2)}}\infty~. 
    \end{align*}
    If $\rela(\vec{y}^{(i-1)}(\vec{x}^{(1)})), \rela(\vec{y}^{(i-1)}(\vec{x}^{(2)})) \in \{0,1\}$, we have the second bullet for $i=i_0$.
    Otherwise, according to \cref{lem:lipschitz-of-the-converter-for-symmetry}, we have 
    \begin{align*}
        \abs{y_2^{(i)}(\vec{x}^{(1)})-y_2^{(i)}(\vec{x}^{(2)})} 
        &\leq 
        2^{7k}\cdot \normtxt{\vec{x}^{(1)}-\vec{x}^{(2)}}\infty + 2^{3n} \cdot \abs{\vec{y}^{(i-1)}(\vec{x}^{(1)}) - \vec{y}^{(i-1)}(\vec{x}^{(2)})}
        \\
        &\leq
        2^{7k}\cdot \normtxt{\vec{x}^{(1)}-\vec{x}^{(2)}}{\infty} + 2^{3n}\cdot (2^{3(i-1)n+7k+\log 2(i-1)}+2^{7k}) \cdot \normtxt{\vec{x}^{(1)}-\vec{x}^{(2)}}{\infty}
        \\
        &\leq
        2^{7k}\cdot \normtxt{\vec{x}^{(1)}-\vec{x}^{(2)}}{\infty} + (2i-1)\cdot 2^{3in+7k}\cdot \normtxt{\vec{x}^{(1)}-\vec{x}^{(2)}}{\infty}
        \\
        &\leq 
        2^{3in+7k + \log 2i}\cdot \normtxt{\vec{x}^{(1)}-\vec{x}^{(2)}}{\infty} ~,
    \end{align*}
    which gives the first bullet for $i=i_0$.

    On the other hand, suppose that the induction hypothesis further gives us the second bullet for $i-1$. 
    That is, for any $j\in [2]$, we have 
    \[
        \vec{y}^{(i-1)}(\vec{x}^{(j)}) \in \set{\left(0,\ 1-y_3^{(i-1)}(\vec{x}^{(j)}),\ y_3^{(i-1)}(\vec{x}^{(j)})\right), ~~ \left(1-y_3^{(i-1)}(\vec{x}^{(j)}),\ 0,\ y_3^{(i-1)}(\vec{x}^{(j)})\right)}~.
    \]
    If $\rela(\vec{y}^{(i-1)}(\vec{x}^{(1)})), \rela(\vec{y}^{(i-1)}(\vec{x}^{(2)})) \in \{0,1\}$, we have the second bullet for $i=i_0$.
    Otherwise, according to \cref{lem:lipschitz-of-the-converter-on-boundaries-for-symmetry}, we have 
    \begin{align*}
        \abs{y_2^{(i)}(\vec{x}^{(1)})-y_2^{(i)}(\vec{x}^{(2)})} 
        &\leq 
        2^{7k}\cdot \normtxt{\vec{x}^{(1)}-\vec{x}^{(2)}}\infty + 2^{2n} \cdot \abs{y_3^{(i-1)}(\vec{x}^{(1)}) - y_3^{(i-1)}(\vec{x}^{(2)})}
        \\
        &\leq
        2^{7k}\cdot \normtxt{\vec{x}^{(1)}-\vec{x}^{(2)}}{\infty} + 2^{2n+7k} \cdot \normtxt{\vec{x}^{(1)}-\vec{x}^{(2)}}{\infty}
        \\
        &\leq 
        2^{2in+7k+1}\cdot \normtxt{\vec{x}^{(1)}-\vec{x}^{(2)}}{\infty} ~, 
    \end{align*}
    which gives the first bullet for $i=i_0$.
\end{proof}

Suppose that \cref{alg:recover-2d-sol-symmetry} fails to give us any solution during the simulation phase, i.e., for any $2\leq i\leq k-1$ and $j\in [3]$, $\mathcal{N}(\vec{y}^{(i)}(\vec{x}^{(j)}))$ is at most bichromatic. 
Since the solution of the \outofk Approximate Symmetric \Sperner instance satisfies $\normtxt{\vec{x}^{(j_1)}-\vec{x}^{(j_2)}}{\infty}\leq 2^{-4kn}$, we have $\abstxt{y_3^{(i)}(\vec{x}^{(j_1)})-y_3^{(i)}(\vec{x}^{(j_2)})} < 2^{-5n}=\eps^5$ for any $2\leq i\leq k$ and any $j_1,j_2\in [3]$, and further that
\begin{itemize}[itemsep=0.2em, topsep=0.5em]
    \item {\bf the intermediate projections are close to each other:} for any $j_1,j_2\in [3]$,$\abstxt{y_2^{(k)}(\vec{x}^{(j_1)})-y_2^{(k)}(\vec{x}^{(j_2)})}<2^{-2n}=\eps^2$ and $\abstxt{y_2^{(i)}(\vec{x}^{(j_1)})-y_2^{(i)}(\vec{x}^{(j_2)})}<2^{-5n}=\eps^5$ for $i<k$; or
    \item {\bf all intermediate projections are on the left/right boundaries:} for any $j\in [3]$, we have either $y_1^{(i)}(\vec{x}^{(j)})=0$ or $y_2^{(i)}(\vec{x}^{(j)})=0$.
\end{itemize}

Next, we prove the same characterization of the intermediate palettes used in \cref{alg:approx-symmetric-sperner-kgeq3} by \cref{lem:characterization-of-intermediate-colors-symmetry}.
The characterization gives equivalence between the set of relevant colors in the palette.
In the first case, where the intermediate projections of the three input vectors are close to each other (and at least one of them is not on the left/right boundaries), the palettes are exactly the same.
In the second case, where the intermediate projections of the input vectors are all on the left/right boundaries, the only two relevant colors for the points, $c^{(i)}(\vec{x}, i^*(\vec{y}^{(i)}(\vec{x})))$ and $c^{(i)}_3(\vec{x})$, are respectively equal. 

\begin{lemma}
    \label{lem:characterization-of-intermediate-colors-symmetry}
    Suppose that $P_{i+1}^{(k-i)}(\vec{x}^{(j)})\leq 0.9$ for any $2\leq i\leq k$ and any $j\in [3]$.
    Suppose that $\mathcal{N}(\vec{y}^{(i)}(\vec{x}^{(j)}))$ is not trichromatic for any $2\leq i<k$ and $j\in [3]$. 
    For any $2\leq i\leq k$, at least one of the following holds:
    \begin{itemize}
        \item The intermediate projections are close to each other, and we have that the corresponding palettes are the same: $\vec{c}^{(i)}(\vec{x}^{(j_1)})=\vec{c}^{(i)}(\vec{x}^{(j_2)})$ for any $j_1,j_2\in [3]$.
        \item All intermediate projections are on the left/right boundaries, and the palettes may be different on an irrelevant color, but we still have that both of the following hold: 
        \begin{itemize}
            \item The color of the first non-zero coordinate $c^{(i)}\left(\vec{x}^{(j)}, i^*(\vec{y}^{(i)}(\vec{x}^{(j)}))\right)$ is the same across all  $j\in [3]$; and
            \item the 3rd color is the same, $c^{(i)}_3(\vec{x}^{(j)})=i+1$ for any $j\in [3]$.
        \end{itemize} 
    \end{itemize}
\end{lemma}

Since the output $\vec{x}^{(1)}, \vec{x}^{(2)}, \vec{x}^{(3)}$ is trichromatic in $\Csym{k}$, we have 
\begin{align}
    \label{eqn:final-trichromatic-symmetry}
    \abs{\set{c^{(k)}(\vec{x}^{(j)},C(\vec{y}^{(k)}(\vec{x}^{(j)})))}_{j\in [3]}}=3~.
\end{align}
We should always have the first bullet of \cref{lem:characterization-of-intermediate-colors-symmetry} for $k$ when each $\mathcal{N}(\vec{y}^{(i)}(\vec{x}^{(j)}))$ is not trichromatic, because otherwise the second bullet of \cref{lem:characterization-of-intermediate-colors-symmetry} violates \cref{eqn:final-trichromatic-symmetry} as $C(\vec{y}^{(k)}(\vec{x}^{(j)}))\in \{i^*(\vec{y}^{(k)}(\vec{x}^{(j)})), 3\}$.
Note that
the first bullet of \cref{lem:characterization-of-intermediate-colors} and \cref{eqn:final-trichromatic} imply that the colors in the base instance $C(\vec{y}^{(k)}(\vec{x}^{(1)}))$, $C(\vec{y}^{(k)}(\vec{x}^{(2)}))$, and $C(\vec{y}^{(k)}(\vec{x}^{(3)}))$ should be distinct. 
We can always guarantee that the tuple $\vec{y}^{(k)}(\vec{x}^{(1)}), \vec{y}^{(k)}(\vec{x}^{(2)}), \vec{y}^{(k)}(\vec{x}^{(3)})$ gives a solution to $C$.
\begin{corollary}
    \label{cor:no-sols-in-sim-implies-final-ones-symmetry}
    Suppose that $P_{i+1}^{(k-i)}(\vec{x}^{(j)})\leq 0.9$ for any $2\leq i\leq k$ and any $j\in [3]$.
    Suppose that $\mathcal{N}(\vec{y}^{(i)}(\vec{x}^{(j)}))$ is not trichromatic for any $2\leq i<k$ and $j\in [3]$. 
    Then, we have $|\settxt{C(\vec{y}^{(k)}(\vec{x}^{(j)})): j\in [3]}|=3$.
\end{corollary}

Before giving the proof of \cref{lem:characterization-of-intermediate-colors-symmetry}, we give an useful lemma here. The lemma states that any point very close (e.g., $<\eps^3$) to a hot point should be either hot or warm. 
This could resolve the problem arisen from the fact that $d^{\alpha}(\cdot, \cdot)$ is only quasimetric and does not enjoy the triangle inequality. 
\begin{lemma}
\label{lem:close-to-hot-points}
    For any $\vec{x}$ that is hot and any $\vec{x'}$ such that $d(\vec{x}, \vec{x'})\leq \eps^3$, $\vec{x'}$ is either hot or warm. 
\end{lemma}
\begin{proof}
    If $\vec{x}, \vec{x'}$ has different colors, both of them are hot. 
    Next, we consider that $\vec{x}, \vec{x'}$ have the same color.
    Let $\vec{y}=\nna(\vec{x})$ denote the nearest neighbor of $\vec{x}$.
    If $x_3,x'_3\geq 0.05$, because of \cref{lem:same-def-for-converted-coordinate} and the fact $d(\cdot,\cdot)$ satisfies the triangle inequality, the lemma holds as 
    \begin{align*}
        d^\alpha(\vec{x'},\vec{y}) = d(\vec{x'},\vec{y}) \leq d(\vec{x}, \vec{x'}) + d(\vec{x}, \vec{y}) = d(\vec{x}, \vec{x'}) + d^{\alpha}(\vec{x}, \vec{y}) < \eps^3 + \eps^2 < 2\eps^2~.
    \end{align*}
    Further, we consider that $x_3,x'_3\leq 0.06$ to finish the proof (because we have $|x_3-x_3'|\leq \eps^3$). 
    Since $\vec{x}$ is hot, we have $d^{\alpha}(\vec{x}, \vec{y})<\eps^2$, and thus
    \begin{align*}
        \alpha(x_3) \cdot |x_2-y_2|, |x_3-y_3| < \eps^2~.
    \end{align*}
    Because $d(\vec{x}, \vec{x'})\leq \eps^3$, we have $|x_2-x'_2|, |x_3-x'_3|\leq \eps^3$, and thus 
    \begin{align*}
        \alpha(x'_3) \cdot |x'_2-y_2| & \leq |\alpha(x_3)-\alpha(x'_3)| + \alpha(x_3) \cdot |x'_2-y_2|\\
        &\leq O(n)\cdot |x_3-x'_3| + \alpha(x_3)\cdot (|x'_2-x_2| + |x_2-y_2|) \tag{\cref{lem:shrinking-factor-basics}}\\
        &\leq O(n)\cdot \eps^3 + \eps^3 + \eps^2 < 2\eps^2~,\\
        |x'_3-y_3| & \leq |x_3-x'_3| + |x_3-y_3| < \eps^3 +\eps^2 < 2\eps^2~.
    \end{align*}
    We have $d^{\alpha}(\vec{x'}, \vec{y})<2\eps^2$. Because $\vec{y}$ (or, an infinite sequence of points that converges to $\vec{y}$) has a different color with $\vec{x},\vec{x'}$, $\vec{x'}$ is hot or warm.
\end{proof}

\begin{proof}[Proof of \cref{lem:characterization-of-intermediate-colors-symmetry}]
    We prove this lemma by induction.
    The base case is when $i=2$. 
    At the beginning of \cref{alg:approx-symmetric-sperner-kgeq3}, we have $\normtxt{\vec{y}^{(2)}(\vec{x}^{(j_1)})-\vec{y}^{(2)}(\vec{x}^{(j_2)})}\infty\leq 2^{O(k)}\cdot 2^{-4kn}<\eps^5$ for any $j_1,j_2\in [3]$ according to the proof of \cref{lem:intermediate-projections-are-close-to-each-other-symmetry} and $\vec{c}^{(2)}(\vec{x}^{(j)})=(1,2,3)$ for any $j\in [3]$. 
    We have the first bullet satisfied for $i=2$. 

    Consider any $i_0\geq 3$. Assume that we have established this lemma for any $i=i_0-1$. 
    Next, we establish this lemma for $i=i_0$.

    First, consider that the first bullet holds for $i-1$. We have $\abstxt{y_2^{(i-1)}(\vec{x}^{(j_1)})-y_2^{(i-1)}(\vec{x}^{(j_2)})}< 2^{-5n}=\eps^5$ for any $j_1,j_2\in [3]$. 
    Hence, $d(\vec{y}^{(i-1)}(\vec{x}^{(j_1)}), \vec{y}^{(i-1)}(\vec{x}^{(j_2)}))\leq \eps^5$ for any $j_1,j_2\in [3]$. 
    We discuss two cases on whether there is a hot point in $\vec{y}^{(i-1)}(\vec{x}^{(1)}), \vec{y}^{(i-1)}(\vec{x}^{(2)})$ and $\vec{y}^{(i-1)}(\vec{x}^{(3)})$. 
    \begin{itemize}[itemsep=0.2em,topsep=0.5em]
        \item W.l.o.g., suppose that $\vec{y}^{(i-1)}(\vec{x}^{(1)})$ is a hot point. 
        According to \cref{lem:lipschitz-of-the-converter-for-symmetry} and \cref{line:y_1,line:y_2,line:y_3}, all the $i$-th intermediate projections are close to each other. 
        According to \cref{lem:close-to-hot-points}, $\vec{y}^{(i-1)}(\vec{x}^{(2)})$ and $\vec{y}^{(i-1)}(\vec{x}^{(2)})$ are also hot or warm.
        Since they are all in a region that is at most bichromatic, 
        the following color set $\{C(\vec{y}^{(i-1)}(\vec{x}^{(j)})), \hat{C}^\alpha_{\sf nn}(\vec{y}^{(i-1)}(\vec{x}^{(j)}))\}$ is then the same across all $j\in [3]$.
        Since $\vec{c}^{(i-1)}(\vec{x}^{(1)})=\vec{c}^{(i-1)}(\vec{x}^{(2)})=\vec{c}^{(i-1)}(\vec{x}^{(3)})$, we then have the same $\vec{c}^{(i)}(\vec{x}^{(j)})$ across all $j\in [3]$, which gives the first bullet of this lemma. 
        \item Otherwise, suppose that $\vec{y}^{(i-1)}(\vec{x}^{(j)})$ is warm or cold for each $j\in [3]$. We have $\rela(\vec{y}^{(i-1)}(\vec{x}^{(j)}))\in \{0,1\}$ for any $j\in [3]$ in this case.
        The $i$-th intermediate projections are on the left/right boundaries, i.e., we have $y_1^{(i)}(\vec{x}^{(j)})=0$ or $y_2^{(i)}(\vec{x}^{(j)})=0$ for each $j\in [3]$.
        Note that there is only one color in $\settxt{C(\vec{y}^{(i-1)}(\vec{x}^{(j)}))}_{j\in [3]}$, because otherwise we clearly have $d^\alpha(\vec{y}^{(i-1)}(\vec{x}^{(j)}),\nna(\vec{y}^{(i-1)}(\vec{x}^{(j)})))<\eps^2$ and all points are hot. 
        Note that \cref{alg:approx-symmetric-sperner-kgeq3} ensures in this scenario that 
        \begin{align*}
            i^*(\vec{y}^{(i)}(\vec{x}^{(j)}))=1 
            &\Leftrightarrow 
            \rela(\vec{y}^{(i-1)}(\vec{x}^{(j)}))=0\\
            &\Leftrightarrow
            C(\vec{y}^{(i-1)}(\vec{x}^{(j)}))<\hat{C}^\alpha_{\sf nn}(\vec{y}^{(i-1)}(\vec{x}^{(j)}))\\
            &\Leftrightarrow
            c^{(i)}(\vec{x}^{(j)},1) = c^{(i-1)}(\vec{x}^{(j)}, C(\vec{y}^{(i-1)}(\vec{x}^{(j)})))
        \end{align*}
        where the second {\em iff} is obtained by \cref{fact:modified-neighboring-color-for-symmetry}.
        We have $c^{(i)}(\vec{x}^{(j)}, i^*(\vec{y}^{(i)}(\vec{x}^{(j)}))) = c^{(i-1)}(\vec{x}^{(j)}, C(\vec{y}^{(i-1)}(\vec{x}^{(j)})))$ for each $j\in [3]$.
        Since $\vec{c}^{(i-1)}(\vec{x}^{(1)})=\vec{c}^{(i-1)}(\vec{x}^{(2)})=\vec{c}^{(i-1)}(\vec{x}^{(3)})$ and $C(\vec{y}^{(i-1)}(\vec{x}^{(1)}))=C(\vec{y}^{(i-1)}(\vec{x}^{(2)}))=C(\vec{y}^{(i-1)}(\vec{x}^{(3)}))$, 
        we have the second bullet for $i$.
    \end{itemize}

    Second, consider that the second bullet holds for $i-1$. 
    Because $|y_3^{(i-1)}(\vec{x}^{(j_1)})-y_3^{(i-1)}(\vec{x}^{(j_2)})|<\eps^{5}$ for any $j_1,j_2\in [3]$, according to the characterization of the temperature on the left/right boundaries (\cref{lem:symmetry-on-converted-coordinates-for-symmetric-sperner}), we have 
    \begin{itemize}[itemsep=0.2em,topsep=0.5em]
        \item for each $j\in [3]$, $\vec{y}^{(i-1)}(\vec{x}^{(j)})$ is hot or warm; or 
        \item for each $j\in [3]$, $\vec{y}^{(i-1)}(\vec{x}^{(j)})$ is warm or cold. 
    \end{itemize}
    Since we have $C(1-z,0,z)=3=C(0,1-z,z)$ or $C(1-z,0,z)=1, C(0,1-z,z)=2$ for any $z\in [0,1]$ (\cref{lem:base-instance-lr-boundaries}), we have the same $c^{(i-1)}(\vec{x}^{(j)},C(\vec{y}^{(i-1)}(\vec{x}^{(j)})))$ across all $j\in [3]$.
    Next, we discuss the above two cases to finish the proof. 
    \begin{itemize}[itemsep=0.2em,topsep=0.5em]
        \item Consider when each $\vec{y}^{(i-1)}(\vec{x}^{(j)})$ is hot or warm. We have $\hat{C}^\alpha_{\sf nn}(\vec{y}^{(i-1)}(\vec{x}^{(j)}))=C^\alpha_{\sf nn}(\vec{y}^{(i-1)}(\vec{x}^{(j)}))$ for each $j\in [3]$. Because we have $C^\alpha_{\sf nn}(1-z,0,z)=1, C^\alpha_{\sf nn}(0,1-z,z)=2$, or $C^\alpha_{\sf nn}(1-z,0,z)=3=C^\alpha_{\sf nn}(0,1-z,z)$ for each $z\in 0.1\pm 2\eps^2$ (\cref{lem:symmetry-on-converted-coordinates-for-symmetric-sperner}), we have the same $c^{(i-1)}(\vec{x}^{(j)},\hat{C}^\alpha_{\sf nn}(\vec{y}^{(i-1)}(\vec{x}^{(j)})))$ across all $j\in [3]$.
        Therefore, the palette $\vec{c}^{(i)}(\vec{x}^{(j)})$ is the same across all $j\in [3]$. 
        Because of \cref{lem:intermediate-projections-are-close-to-each-other-symmetry} and the Lipschitzness of the coordinate converter on the left/right boundaries \cref{lem:lipschitz-of-the-converter-on-boundaries-for-symmetry}, the $i$-th intermediate projections, $\vec{y}^{(i)}(\vec{x}^{(1)}), \vec{y}^{(i)}(\vec{x}^{(2)})$ and $\vec{y}^{(i)}(\vec{x}^{(3)})$, are close to each other. 
        Hence, we prove the first bullet of this lemma for $i$. 
        \item Consider when each $\vec{y}^{(i-1)}(\vec{x}^{(j)})$ is warm or cold. We have $\rela(\vec{y}^{(i-1)}(\vec{x}^{(j)}))\in \{0,1\}$ and all $i$-th intermediate projections are on the left/right boundaries. 
        Note that there is only one color in $\settxt{C(\vec{y}^{(i-1)}(\vec{x}^{(j)}))}_{j\in [3]}$, because otherwise we clearly have $d^\alpha(\vec{y}^{(i-1)}(\vec{x}^{(j)}),\nn(\vec{y}^{(i-1)}(\vec{x}^{(j)})))<\eps^2$ and all points are hot. 
        According to our earlier discussions, \cref{alg:approx-symmetric-sperner-kgeq3} ensures in this scenario that $c^{(i)}(\vec{x}^{(j)}, i^*(\vec{y}^{(i)}(\vec{x}^{(j)}))) = c^{(i-1)}(\vec{x}^{(j)}, C(\vec{y}^{(i-1)}(\vec{x}^{(j)})))$.
        Therefore, we have the same $c^{(i)}(\vec{x}^{(j)}, i^*(\vec{y}^{(i)}(\vec{x}^{(j)})))$ across all $j\in [3]$, and thus the second bullet of this lemma holds. \qedhere
    \end{itemize}
\end{proof}

\paragraph{Case 2: the output does not satisfy~\cref{eqn:assumption-for-simple-proof-of-recovery-symmetry}.}
Next, we complete our second step by showing how to prove \cref{lem:poly-time-recovery-symmetry} without the property (\cref{eqn:assumption-for-simple-proof-of-recovery-symmetry}). 
Suppose $\theta^*$ is the minimum threshold $\theta\geq 2$ such that $P_{i+1}^{(k-i)}(\vec{x}^{(j)})\leq 0.9$ for any $\theta\leq i\leq k$ and $j\in[3]$. 
Such threshold always exists because otherwise we have $x^{(1)}_{k+1}, x^{(2)}_{k+1}, x^{(3)}_{k+1}\geq 0.8$. 
This implies $y_3^{(k)}(\vec{x}^{(j)})>0.8$ for any $j\in [3]$, and according to our construction of the base instance (\cref{alg:base-instance}), we have $C^{(k)}(\vec{x}^{(j)})=c^{(k)}_3(\vec{x}^{(j)}) = k+1$ for any $j\in[3]$, violating the assumption that $\vec{x}^{(1)}, \vec{x}^{(2)}, \vec{x}^{(3)}$ form a solution for the \outofk Approximate Symmetric \Sperner problem. 
When $\theta^*=2$, it is equivalent with the special case satisfying \cref{eqn:assumption-for-simple-proof-of-recovery-symmetry} and we have proved \cref{lem:poly-time-recovery-symmetry} for this case.
On the other hand, if $\theta^*>2$, we have $P^{(k-\theta^*+1)}_{\theta^*}(\vec{x}^{(j)})\geq 0.8$ for any $j\in [3]$ because of the Lipschitzness of the projection step \cref{lem:lipschitz-of-projection} and that $\normtxt{\vec{x}^{(i)}-\vec{x}^{(j)}}{\infty}\leq 2^{-4kn}$ for any $i,j\in [3]$. 
This means that we have $C(\vec{y}^{(\theta^*-1)}(\vec{x}^{(j)}))=3$.
And because of \cref{lem:trivial-converter-for-symmetry}, we have $\rela(\vec{y}^{(\theta^*-1)}(\vec{x}^{(j)}))=1$ and $y_1^{(\theta^*)}(\vec{x}^{(j)})=0$ for any $j\in [3]$. 
Further, we have for each $j\in [3]$,
\begin{align}
    \label{eqn:y-theta-star}
    \vec{y}^{(\theta^*)}(\vec{x}^{(j)})=\left(0,\ 1-P^{(k-\theta^*)}_{\theta^*+1}(\vec{x}^{(j)}),\ P^{(k-\theta^*)}_{\theta^*+1}(\vec{x}^{(j)})\right)~.
\end{align}
Therefore, running \cref{alg:approx-symmetric-sperner-kgeq3} on instance $C^{(k)}$ for $\vec{x}^{(1)},\vec{x}^{(2)}, \vec{x}^{(3)}$ is equivalent to running \cref{alg:approx-symmetric-sperner-kgeq3} on instance $C^{(k')}$ with $k'=k-\theta^*+2$ for $\vec{\hat{x}}^{(1)},\vec{\hat{x}}^{(2)},\vec{\hat{x}}^{(3)}\in \Delta^{k'}$ such that
\begin{align*}
    \forall i\in [k'+1], j\in [3], \quad \hat{x}_i^{(j)} = \begin{cases}
        0 & \text{if $i=1$,}\\
        1-\sum_{i'=\theta^*+1}^{k+1} x^{(j)}_{i'} & \text{if $i=2$,}\\
        x^{(j)}_{i+\theta^*-2} & \text{if $i>2$.}
    \end{cases}
\end{align*}
This is because $\vec{y}^{(2)}(\vec{\hat{x}}^{(j)})$ equals the RHS of \cref{eqn:y-theta-star}.
Because this new instance has a smaller number of dimensions and satisfies the condition of our special cases (\cref{eqn:assumption-for-simple-proof-of-recovery-symmetry}), we complete the proof for \cref{lem:poly-time-recovery-symmetry}.